\chapter{面向受限条件的协同增量语义 HARQ 可靠传输机制}
\label{chap:semantic}

\section{引言}

沿“感知—接入—调度—传输”的端到端数据通路，第二章在退化遥测条件下构建了可信的本地多跳网络状态感知，第~\ref{chap:handover}~章保障了海量终端接入的连续性与公平性，第~\ref{chap:resource}~章在受限星上资源下实现了面向用户级服务质量（Quality of Service，QoS）的跨层调度。然而，可信感知、公平接入与高效调度最终仍须落实为业务数据在受限链路上的可靠送达：在星上算力与链路带宽双重受限的条件下，如何以尽可能低的开销将任务相关信息可靠传输至接收端，决定了端到端 QoS 保障能否闭环。由此，本章从“调度”推进至“传输”，研究纯分布式管控范式下面向受限条件的语义可靠传输机制。

在第六代移动通信（6G）空天地一体化网络中，大规模 LEO 星座承载的业务流量中相当一部分为面向下游视觉任务的高分辨率遥感图像~\cite{xiao-SAGINSurvey-JSAC2024}，其在地面侧主要用于场景分类~\cite{neware-ClassificationSurvey-ICECA2018}、语义分割~\cite{meng-SwinTransformer-GRS2022}等智能解译任务，而非像素级精确重建。此类图像数据量大，叠加星地/星间链路带宽受限与误码起伏，使其高效可靠传输成为突出难题。由于图像交付的根本目的在于支撑下游任务执行而非逐比特恢复，传统比特级传输会将稀缺的通信资源耗费在对任务无贡献的视觉冗余上，传输效率低下~\cite{Shao-TaskOrientedCommunication-JSAC2022}。视觉语义通信~\cite{peng-RobustImageSemanticCommunication-JSAC2025,Park-TransmitWhatYouNeed-JSAC2025}将传输范式由比特恢复转向任务相关语义提取，仅传输对下游任务有用的语义特征，可在维持任务性能的同时显著降低带宽消耗与端到端时延，为受限星上链路提供了新的可靠传输思路。其中，联合信源信道编码（JSCC）通过端到端学习建立源数据与信道符号之间的鲁棒映射~\cite{Bourtsoulatze-DeepJSCC-ICASSP19}，构成语义图像传输的基础骨干。

然而，现有语义通信框架大多是任务导向的~\cite{xie-TaskOrientedMultiUserSemanticCommunication-JSAC2022,zhang-MultiTaskSemanticCommunication-TCOM2024}，其编解码器针对特定下游目标联合优化：编码端利用任务先验保留任务相关特征、舍弃语义无关内容，从而形成高度任务相关的语义表征。这种语义编码与任务目标的紧耦合，与实际星上图像交付系统并不相容。在 LEO 卫星图像交付中，星上发送端通常仅负责图像获取与转发，下游任务则在接收端执行，且同一图像可能支撑多种异构任务~\cite{xiao-SurveyFoundationModels-GRS2025}；若为每类任务分别训练并维护专用编解码器，将带来过高的复杂度与较差的可扩展性。更进一步，高性能语义传输所依赖的大参数量 Transformer 编码与生成式复原模型，其计算、存储与时延开销过高，难以部署于受限星上载荷。在发送端任务无关、接收端任务感知的非对称架构下，重传的角色因而被重新定义~\cite{zhang-TaskUnawareSemanticCommunication-JSAC2023}：发送端无法判断哪些语义分量对未知下游任务尚不充分，须由接收端以紧凑形式回指其所缺失的任务相关信息。HARQ 经由确认/否认（ACK/NACK）反向控制通路~\cite{Sheng-SCCoopPerception-TCCN26,zhou-ABRControl-HARQ-SemanticCommunication-OJCOMS2022}，恰好为承载此类语义反馈提供了天然接口，从而引出由接收端反馈驱动、跨轮逐步精化初始任务无关表征的增量传输框架。

围绕语义 HARQ，已有研究在细粒度错误定位与反馈驱动自适应方面取得进展。早期工作以句级语义相似度替代比特级校验，并按信道条件与信源复杂度自适应调整语义码率~\cite{jiang-DeepSourceChannelCoding-TCOM2022,zhou-ABRControl-HARQ-SemanticCommunication-OJCOMS2022}；其后，特征级失真评估~\cite{hu-SemHARQ-JSAC25}与基于语义码本的细粒度错误识别~\cite{liang-SemanticCodebookHARQ-TCOMM25}进一步将重传触发由比特正确性推进至语义可靠性；近期研究则引入生成式隐空间模型实现多轮语义合并~\cite{han-GenerativeSemanticHARQ-arXiv2026}，以及在联合信源信道编码框架内基于反馈的重传决策与再编码~\cite{li-Boyuan-JS3C-GLOBECOM2025}。然而，上述方案在任务无关的星上图像交付场景下仍存在若干局限。其一，多数机制以重建或感知质量为导向，刻画的是整体输出质量而非单个语义分量对任务结果的贡献，难以定位最能修复当前任务缺口的重传分量。其二，语义表征多以任务耦合方式学习，要求发送端预知下游任务，限制了其在任务无关场景下的适用性；即便在任务无关发送端设定下，接收端也未显式量化任务执行所缺失的具体语义分量，发送端亦未在严格反馈预算下逐轮精化语义表征。其三，反馈信令多局限于 ACK/NACK 或标量估计等粗粒度指示，不足以刻画具体的语义缺陷；多轮接收端则多依赖特征平均或预设合并规则，未显式建模任务相关语义在多轮传输中的演化与互补。因此，现有语义 HARQ 仍缺乏将接收端任务缺口定位、紧凑反馈信令、反馈条件化语义增量生成与跨轮状态感知累积相贯通的闭环精化机制。

针对上述问题，本章面向任务无关发送端与任务感知接收端的非对称 LEO 图像交付场景，提出协同增量语义混合自动重传请求（Coordinated Incremental Semantic Hybrid Automatic Repeat Request，CIS-HARQ）机制，将重传由对既有语义的重复发送重构为反馈条件化的选择性精化过程。具体而言，接收端在严格反馈预算下，将由任务输出裕量梯度刻画的任务相关语义缺口压缩为紧凑的 Top-$K$ 硬掩码并回传；发送端据此结合 HARQ 轮次嵌入，仅对被请求分量生成选择性语义增量；接收端再以门控融合与掩码引导的语义累积仅更新被请求分量，在保留有用历史语义、抑制含噪与冗余增量的同时，于前向链路、反馈链路与模型复杂度三重约束下逐步恢复任务相关语义。由此，轻量化收发机得以经多轮收发协同逐步逼近高复杂度模型的任务性能。

本章的主要贡献总结如下：
\begin{itemize}
\item 面向任务无关发送端与任务感知接收端的非对称图像交付，提出闭环语义 HARQ 框架 CIS-HARQ。该框架在不向发送端暴露下游任务模型的前提下，经由接收端反馈实现任务自适应重传，使发送端无需为每类下游任务维护专用编解码器。

\item 针对反馈生成，提出任务驱动的语义缺口辨识机制，将接收端的任务不确定性（由任务输出裕量梯度与特征激活的联合归因刻画）映射为紧凑的分块级 Top-$K$ 硬掩码，使反馈在严格预算下仍与低速 HARQ 控制信令相兼容。

\item 针对增量编码，设计反馈条件化的语义精化编码器，将硬掩码与 HARQ 轮次嵌入转化为选择性语义增量，使发送端仅精化接收端所请求的分量而非重传完整表征，并以速率约束抑制增量信息密度，控制附加 HARQ 负载。

\item 针对接收端合并，提出掩码引导的语义累积机制，跨轮选择性更新被请求分量，保护有用历史语义免受含噪增量干扰，并支撑多轮互补精化。

\item 在 ImageNet-100、CIFAR-100 与 VeRi-776 数据集上的实验表明，CIS-HARQ 在中低信道带宽比下通过反馈引导精化显著提升下游任务精度，紧带宽下经三轮增量将分类精度由同带宽基线的 $35.7\%$ 提升至 $65.0\%$，以远低于生成式与 Transformer 基线的参数规模与推理时延部分弥合其性能差距，并在不同数据分布的数据集与多任务接收端上保持一致的精化增益。
\end{itemize}

\section{本章小结}
